% Nouveau EGA V (cote 25): archive cote 25, site page 117 (PDF page 118).
% typescript. Draft compared: none.
% AI-transcribed by Claude Opus 5.5 (claude-opus-5-5, adaptive thinking, effort high), Message Batches API, batch ega25b (msgbatch_01A7n5YuXubUPNK1zsepi8RG), 2026-11-13; not yet checked by eye.
\scanpage{25/117}
\mspage{20}

de diviseurs de $X/S$ paramétrée par $T$, on appelle ``composante fixe'' de la
famille toute composante irréductible de dodimension 1 de l'ensemble $Z$ de\unsure{s}
points fixes de la famille; on dit que la famille est ``sans composante
fixe'' si elle est sans composante fixe, i.e.\ si $\operatorname{codim}(Z,X) \hand{$\geq$} 2$. Cette
terminologie s'étend aussitôt au cas où $S$ est quelconque, en regardant
fibre par fibre. Le fait d'être sans composante fixe est évidemment stabl\unsure{e}
par changement de base.

\uline{Proposition}\over{2}{\illegible{1}} \hand{/}Supposons $X \longrightarrow S$ plat loc de présentation finie, à fibres
$(S_2)$ , et soit $D$ un syst\xout{e}ême linéaire de diviseurs \uline{sans composante fixe}
sur $X/S$, paramétrée par $\overset{\hand{$\vee$}}{P}$ . Alors $D$ est uniquement déterminé par la
connaissance du morphisme $f: X - Z \longrightarrow P$ correspondant ($Z$ = ensemble
des points fixes), et même par la connaissance de la classe de $f$ comme
\hand{``}\deleted{application rationelle}\inserted{pseudo-morphisme} \xout{aux dessus de S}\inserted[typed]{rel à S}\hand{''} ; $X-Z$ est le \xout{plus grand ouvert}\inserted[typed]{domain\xout{x}e}
de définition de ladite classe.

Pour la notion et le sorite des ``\deleted{applications rationelles}\inserted{pseudo-morphismes}
rel à S'' voir \deleted{SGAD XVIII n°1}\inserted{\S 20.10.} , \deleted{qu'on pourrait éventuellement rajouter
au dernier n° du par.11 sur la densité schématique.} Il faut prouver que.
si \xout{U est un ouvert de X schématiquement relativement dense} $D'$ est un\deleted{e}
autre syst lin de diviseurs sans comp fixe paramétrée par $\overset{\hand{$\vee$}}{P}$ , \xout{xxx}
définissant $f': X-Z' \longrightarrow P$, et si $f$ et $f'$ coincident dans un ouvert $U$
$\hand{$\subset$} (X-Z)\hand{$\cap$}(X-Z')$ schématiquement dense relativement à $S$, alors $D=D'$ .
En fait, comme $P$ est séparé sur $S$, on peut prendre $U=(X-Z)\hand{$\cap$}(X-Z') =$
\xout{X - (Z u Z') = d'autre part} $X-Z''$ , où $Z'' = Z\hand{$\cup$}Z'$ . Comme $Z''$ est de codim$\hand{$\geq$}2$
sur chaque fibre et que $X$ est à fibres $(S_2)$, il s'ensuit que pour tout
$x\hand{$\in$} Z''$, la fibre $X_s$ est de profondeur $\hand{$\geq$}2$ en $x$. On va en conclure (utili-%
sant de plus bien sûr le fait que $X$ est plat loc de présentation finie
sur $S$) que tout diviseur \deleted{de Cartier positif} sur $X$ (pas nécessairement \margin{\unsure{cf}\ \S 20.11.}
transversal aux fibres) est connu quand on connait sa restriction à $X-Z''$

\mspage{117}
